\documentclass[aspectratio=169,11pt]{beamer}
\usepackage[utf8]{inputenc}
\usepackage[T1]{fontenc}
\usepackage[french]{babel}
\usepackage{helvet}
\renewcommand{\familydefault}{\sfdefault}
\usepackage{hyperref}
\definecolor{bleu}{HTML}{0000FF}
\definecolor{fond}{HTML}{F4F4EF}
\setbeamercolor{background canvas}{bg=fond}
\setbeamercolor{structure}{fg=bleu}
\setbeamercolor{title}{fg=bleu}
\setbeamercolor{frametitle}{fg=bleu,bg=fond}
\setbeamercolor{normal text}{fg=black}
\setbeamertemplate{navigation symbols}{}
\setbeamertemplate{itemize item}{\textcolor{bleu}{\tiny$\blacksquare$}}
\setbeamertemplate{footline}{\hfill\footnotesize\insertframenumber\hspace{1em}\vskip6pt}
\hypersetup{colorlinks=true,urlcolor=bleu,linkcolor=bleu}
\title{Critiques et débats}
\subtitle{Organisation de coopération et de développement économiques}
\author{Veille institutionnelle}
\date{2026-11-03}

\begin{document}
\begin{frame}\titlepage\end{frame}

\begin{frame}{Critiques et débats (1/2)}
\begin{itemize}
\item Club de pays riches : représentativité limitée des pays en développement
\item Impôt minimum mondial : portée, exceptions, accord de 2025 entre le G7 et les États-Unis (à vérifier)
\item Enquête PISA : méthodes de comparaison et usage des classements
\item Recommandations de réformes structurelles jugées trop libérales par certains
\end{itemize}
\end{frame}

\begin{frame}{Critiques et débats (2/2)}
\begin{itemize}
\item Rôle dans les listes de juridictions non coopératives en fiscalité
\end{itemize}
\end{frame}

\begin{frame}{Critiques et débats : ressources à suivre}
\small
Veille hebdomadaire.
\begin{itemize}
\item \href{https://www.oecd.org/en/topics/policy-areas/tax.html}{Fiscalité et BEPS}\\{\footnotesize Page thématique, publication : hebdomadaire}
\item \href{https://www.oecd.org/en/about/programmes/pisa.html}{PISA}\\{\footnotesize Enquête, publication : annuelle}
\end{itemize}
\end{frame}

\end{document}
